\section{A brief introduction to Lean syntax}

Lean is both an interactive theorem prover and a functional programming language.
Definitions, statements, and proofs are written in a syntax that is
subsequently verified by the Lean kernel. In this section, we introduce several
foundational pieces of Lean syntax that appear frequently in later chapters.

\subsection{Basic declarations}

A typical Lean file consists of a series of declarations. Beyond the foundational
\lstinline|def|, \lstinline|theorem|, and \lstinline|axiom|, Lean provides
powerful organizational tools such as \lstinline|abbrev|, \lstinline|structure|,
\lstinline|class|, and \lstinline|instance|.

The keyword \lstinline|def| is used to define new data, functions, or values.
For example, the following code defines the successor of a natural number:

\begin{lstlisting}[language=Lean]
def succ (n : Nat) : Nat := n + 1
\end{lstlisting}

Here, \lstinline|succ| is the name; \lstinline|(n : Nat)| declares an explicit input \lstinline|n| of type \lstinline|Nat|; the final \lstinline|: Nat| gives the return type; and \lstinline|:=| introduces the defining expression on the right.

Closely related is the keyword \lstinline|abbrev|. It also creates a definition, but instructs Lean's elaborator to treat the new name as fully transparent during unification. It is frequently used for type aliases:

\begin{lstlisting}[language=Lean]
abbrev Vector3 := Float × Float × Float
\end{lstlisting}

Lean then automatically unfolds \lstinline|Vector3| into a triplet of floats during type checking, avoiding the strict type-mismatch errors that a standard \lstinline|def| might cause.

The keyword \lstinline|theorem| is reserved for stating and proving propositions. A theorem requires a name, a logical statement, and a formal proof. For example:

\begin{lstlisting}[language=Lean]
theorem add_zero (a : Nat) : a + 0 = a := by
  rw [Nat.add_zero]
\end{lstlisting}

This theorem asserts that for every natural number \lstinline|a|, the equality \lstinline|a + 0 = a| holds. The portion before \lstinline|:=| is the mathematical statement, while the portion after \lstinline|:= by| enters ``tactic mode'' to construct the proof.

The keyword \lstinline|axiom| introduces a statement assumed to be true without providing a proof:

\begin{lstlisting}[language=Lean]
axiom my_axiom : P
\end{lstlisting}

Axioms instruct Lean's kernel to accept \lstinline|P| unconditionally and must be used sparingly to avoid compromising the soundness of the logical system.

To bundle multiple pieces of data together into a single mathematical object or record, Lean uses the \lstinline|structure| keyword.

\begin{lstlisting}[language=Lean]
structure Point where
  x : Float
  y : Float
\end{lstlisting}

This declaration creates a new type called \lstinline|Point|. It automatically generates a
constructor to build points and projection functions (\lstinline|Point.x| and \lstinline|Point.y|) to access the underlying fields.

Finally, Lean handles algebraic hierarchies and overloading via the
\lstinline|class| and \lstinline|instance| keywords. A \lstinline|class| is
essentially a \lstinline|structure| that registers itself with Lean's typeclass
synthesis mechanism (which resolves the square brackets \lstinline|[]| discussed in
Section~\ref{sec:brackets}).

\begin{lstlisting}[language=Lean]
class Add (α : Type) where
  add : α → α → α
\end{lstlisting}

This defines an interface requiring an addition operation for some type \lstinline|α|. To fulfill this requirement for a specific type, we use the \lstinline|instance| keyword to provide the concrete implementation:

\begin{lstlisting}[language=Lean]
instance : Add Nat where
  add x y := x + y
\end{lstlisting}

Once this instance is declared, any function requiring an \lstinline|Add| instance for \lstinline|Nat| will automatically find and use this implementation.

\subsection{The meaning of different brackets}
\label{sec:brackets}

Lean utilizes different types of brackets to distinguish how arguments should be provided or inferred by the system.

Parentheses \lstinline|()| denote \emph{explicit arguments}. These are arguments that the user must manually provide when calling a function or applying a theorem.

\begin{lstlisting}[language=Lean]
def double (n : Nat) : Nat := n + n
\end{lstlisting}

In this example, \lstinline|(n : Nat)| signifies that \lstinline|n| is an explicit argument. Whenever \lstinline|double| is used, the user must explicitly pass a natural number.

Curly brackets \lstinline|{}| denote \emph{implicit arguments}. Lean attempts to infer these arguments automatically from the surrounding context using unification.

\begin{lstlisting}[language=Lean]
def identity {α : Type} (x : α) : α := x
\end{lstlisting}

Here, \lstinline|{α : Type}| indicates that \lstinline|α| is an implicit type parameter. If we apply \lstinline|identity| to a specific natural number, Lean automatically infers that \lstinline|α| must be \lstinline|Nat|, saving the user from writing redundant types.

Square brackets \lstinline|[]| are designated for \emph{typeclass arguments}. While also inferred automatically, they rely on Lean's typeclass synthesis mechanism rather than simple unification.

\begin{lstlisting}[language=Lean]
def addSelf [Add α] (x : α) : α := x + x
\end{lstlisting}

The expression \lstinline|[Add α]| specifies that Lean must find a registered addition structure (a typeclass instance) for the type \lstinline|α|. Lean searches its database to satisfy this requirement automatically when the function is invoked.

\subsection{Attribute tags}
\label{sec: tag_ext_simp}

Lean supports attribute tags, which append special metadata to declarations
to guide automated tools. An attribute is declared directly above a
definition, theorem, or structure using square brackets preceded by an
\lstinline|@| symbol.

The \lstinline|simp| tactic invokes Lean's simplification engine, which applies
a large library of pre-registered lemmas to reduce a goal, often solving
straightforward goals automatically.

A widely used attribute is \lstinline|@[simp]|. Tagging a theorem with \lstinline|@[simp]| adds it to this simplification database, allowing the \lstinline|simp| tactic to apply it autonomously in future proofs.

\begin{lstlisting}[language=Lean]
@[simp]
theorem add_zero (a : Nat) : a + 0 = a := by
  rw [Nat.add_zero]
\end{lstlisting}

Once marked, whenever the simplifier encounters an expression matching \lstinline|a + 0|, it will automatically rewrite it to \lstinline|a| without requiring explicit user instruction.

Another essential attribute is \lstinline|@[ext]|, which establishes extensionality principles. Extensionality asserts that two objects are equal whenever all their corresponding components are equal. Applying \lstinline|@[ext]| directly to a \lstinline|structure| automatically generates the corresponding extensionality lemma.

\begin{lstlisting}[language=Lean]
@[ext]
structure Point where
  x : Float
  y : Float
\end{lstlisting}

Tagging the \lstinline|Point| structure with \lstinline|@[ext]| instructs Lean's metaprogramming framework to generate a theorem behind the scenes. Written manually, it would look like this:

\begin{lstlisting}[language=Lean]
theorem Point.ext {p q : Point} (h₁ : p.x = q.x) (h₂ : p.y = q.y) : p = q
\end{lstlisting}

The \lstinline|@[ext]| tag registers this logic so that invoking the \lstinline|ext| tactic on a goal of \lstinline|p = q| will automatically apply \lstinline|Point.ext|, splitting the goal into the simpler subgoals \lstinline|p.x = q.x| and \lstinline|p.y = q.y|.

Attributes are indispensable for scaling mathematical libraries.
Rather than manually proving and referencing hundreds of structural lemmas
by name, developers register them with tags like \lstinline|@[simp]| and
\lstinline|@[ext]|, enabling Lean's automation to carry the burden in complex
proofs.

\subsection{Namespaces}
\label{sec: namespace}

% \subsection{Namespaces and the \lstinline|open| Keyword}

As mathematical libraries grow, multiple theorems or definitions with similar names become common. Lean uses namespaces to organize declarations and prevent such naming conflicts. When working extensively within a specific context, such as multivariate power series, repeatedly typing the full prefix \lstinline|MvPowerSeries| becomes tedious and clutters the code. Enclosing our code in a \lstinline|namespace| block lets Lean resolve local names automatically, so we may omit the prefix for any definition or theorem belonging to that namespace.

For instance, consider the following explicit definition of a variable $X$:

\begin{lstlisting}[language=Lean]
def MvPowerSeries.X (s : σ) : MvPowerSeries σ R :=
    (MvPowerSeries.monomial (Finsupp.single s 1)) 1
\end{lstlisting}

If we declare that we are operating inside the \lstinline|MvPowerSeries| namespace, the code becomes significantly cleaner. We can drop the prefix from both the new definition's name (\lstinline|X|) and any internal references to existing functions in that namespace (such as \lstinline|monomial|):

\begin{lstlisting}[language=Lean]
namespace MvPowerSeries

def X (s : σ) : MvPowerSeries σ R :=
    monomial (Finsupp.single s 1) 1

end MvPowerSeries
\end{lstlisting}

Once outside the \lstinline|end MvPowerSeries| boundary, a user would once again need to refer to this definition by its fully qualified name: \lstinline|MvPowerSeries.X|.

A complementary tool for managing scopes is the \lstinline|open| keyword.
While \lstinline|namespace| is typically used when \textit{defining} new
objects within a specific theory, \lstinline|open| is used when we simply
want to \textit{access} existing results from another namespace without
typing its prefix. For example, suppose we have a definition that relies on
\lstinline|Finsupp.single|. If we want to use results from the
finitely supported functions library frequently, we can open the
\lstinline|Finsupp| namespace directly:

\begin{lstlisting}[language=Lean]
open Finsupp
namespace MvPowerSeries

def X (s : σ) : MvPowerSeries σ R :=
    monomial (single s 1) 1

end MvPowerSeries
\end{lstlisting}

By writing \lstinline|open Finsupp|, we instruct Lean to look inside the \lstinline|Finsupp| namespace during name resolution, allowing us to reduce \lstinline|Finsupp.single| to just \lstinline|single|.

% \begin{remark}
% The specific functions presented in this example, such as
% \lstinline|monomial| and \lstinline|single|, are utilized here purely
% to illustrate Lean's syntax and scoping mechanisms. We will formally discuss
% the detailed mathematical API for each of these components in Section~\ref{sec:mvpowerseries}.
% \end{remark}


% As mathematical libraries grow, it is common to have multiple theorems or
% definitions with similar names. Lean uses namespaces to organize declarations
% and prevent these naming conflicts. When working extensively within a
% specific mathematical context, such as multivariate power series,
% repeatedly typing the full prefix \lstinline|MvPowerSeries| can become
% tedious and clutter the code. By encapsulating our code within a
% \lstinline|namespace| block, Lean automatically resolves local
% names, allowing us to omit the prefix for any definitions or theorems
% belonging to that namespace.

% For instance, consider the following explicit definition of a variable $X$:

% \begin{lstlisting}[language=Lean]
% def MvPowerSeries.X (s : σ) : MvPowerSeries σ R :=
%     (MvPowerSeries.monomial (Finsupp.single s 1)) 1
% \end{lstlisting}

% If we declare that we are operating inside the \lstinline|MvPowerSeries| namespace, the code becomes significantly cleaner. We can drop the prefix from both the new definition's name (\lstinline|X|) and any internal references to existing functions in that namespace (such as \lstinline|monomial|):

% \begin{lstlisting}[language=Lean]
% namespace MvPowerSeries

% def X (s : σ) : MvPowerSeries σ R :=
%     monomial (Finsupp.single s 1) 1

% end MvPowerSeries
% \end{lstlisting}

% Once outside the \lstinline|end MvPowerSeries| boundary, a user would once again need to refer to this definition by its full, fully qualified name: \lstinline|MvPowerSeries.X|. This mechanism keeps the global namespace clean while ensuring that local development remains readable and concise.